\medskip

Une entreprise a créé une Foire Aux Questions (\og FAQ \fg) sur son site internet.



On étudie le nombre de questions qui y sont posées chaque mois.

\bigskip

\textbf{Partie A : Première modélisation}

\medskip

Dans cette partie, on admet que, chaque mois :

\begin{itemize}
\item[$\bullet~~$]90\,\% des questions déjà posées le mois précédent sont conservées sur la FAQ ;
\item[$\bullet~~$]130 nouvelles questions sont ajoutées à la FAQ.
\end{itemize}

Au cours du premier mois, $300$ questions ont été posées.



Pour estimer le nombre de questions, en centaines, présentes sur la FAQ le $n$-ième
mois, on modélise la situation ci-dessus à l'aide de la suite $\left(u_n\right)$ définie par : $u_1 =3$ et, pour tout entier naturel $n \geqslant 1$, $u_{n+1} = 0,9u_n + 1,3$.


\begin{enumerate}
\item $u_2 = 0,9\times u_1 +1,3 = 0,9\times 3 + 1,3= 4$ et 
$u_3 =0,9\times u_2 +1,3 = 0,9\times 4 + 1,3= 4,9$


On peut estimer à 400 le nombre de questions le 2\up{e} mois, et à 490 le 3\up{e} mois.

\item On va montrer par récurrence la propriété  $u_n = 13 - \dfrac{100}{9} \times 0,9^n$ pour tout $n\geqslant1$.

\begin{list}{\textbullet}{}
\item \textbf{Initialisation}

Pour $n=1$, on a $u_1=3$ et $13 -\dfrac{100}{9}\times 0,9^1=13-\dfrac{90}{9}=3$.

La propriété est vérifiée au rang 1.
\item \textbf{Hérédité}

On suppose la propriété vraie au rang $n$ avec $n\geqslant 1$; autrement dit $u_n = 13 - \dfrac{100}{9} \times 0,9^n$ ; c'est l'hypothèse de récurrence.

$\begin{gathered}u_{n+1} = 0,9u_n+1,3 = 0,9 \left (13 - \dfrac{100}{9} \times 0,9^n\right ) +1,3
= 0,9\times 13 - \dfrac{100}{9}\times 0,9^{n+1}+0,1\times 13\\
\phantom{u_{n+1}}
= 13- \dfrac{100}{9}\times 0,9^{n+1}\end{gathered}$

Donc la propriété est vraie au rang $n+1$.
\item \textbf{Conclusion}

La propriété est vraie au rang 1 et elle est héréditaire pour tout $n\geqslant 1$; d'après le principe de récurrence, elle est vraie pour tout $n\geqslant 1$.
\end{list}

Donc pour tout $n\geqslant 1$, on a: $u_n = 13 - \dfrac{100}{9} \times 0,9^n$.

\item  
Pour tout $n\geqslant 1$, on a:

$\begin{gathered}u_{n+1}-u_n
=\left (13-\dfrac{100}{9}\times 0,9^{n+1}\right ) - \left (13-\dfrac{100}{9}\times 0,9^{n}\right ) 
= 13-\dfrac{100}{9}\times 0,9^{n+1} - 13+\dfrac{100}{9}\times 0,9^{n}\\
\phantom{u_{n+1}-u_n}
= \dfrac{100}{9}\times 0,9^{n}\left (1-0,9\right )
= \dfrac{100}{9}\times 0,9^{n} \times 0,1>0\end{gathered}$

donc la suite $(u_n)_{n\in\N*}$ est croissante.

\item On considère le programme ci-contre, écrit en langage Python.

\begin{minipage}{0.72\linewidth}



Ce programme renvoie la plus petite valeur $n$ telle que $u_n>\texttt{p}$; donc \texttt{seuil(8.5)} renvoie la plus petite valeur $n$ telle que $u_n>8,5$. On résout cette inéquation.

$\begin{gathered}u_n>8,5
\iff
13 -\dfrac{100}{9}\times 0,9^n > 8,5
\iff
4,5 > \dfrac{100}{9}\times 0,9^n\\
\phantom{u_n>8,5}
\iff
\dfrac{4,5\times 9}{100}> 0,9^n
\iff
\ln(0,405) > \ln \left (0,9^n\right )\\
\phantom{u_n>8,5}
\iff
\ln(0,405) >n \times \ln  (0,9)
\iff
\dfrac{\ln(0,405)}{\ln(0,9)} <n\end{gathered}$
\end{minipage}\hfill
\begin{minipage}{0.25\linewidth}
\renewcommand\arraystretch{0.9}
\begin{tabularx}{\linewidth}{|X|} \hline
def seuil(p) :\\
\qquad n=1\\
\qquad u=3\\
\qquad while u<=p :\\
\qquad  \qquad n=n+1\\
\qquad \qquad u=0.9*u+1.3 \\
\qquad return n\\ \hline
\end{tabularx}
\end{minipage}

$\dfrac{\ln(0,405)}{\ln(0,9)} \approx 8,6$ donc la valeur renvoyée par \texttt{seuil(8.5)} est 9.
\end{enumerate}

\bigskip

\textbf{Partie B : Une autre modélisation}

\medskip

Dans cette partie, on considère une seconde modélisation à l'aide d'une nouvelle
suite $\left(v_n\right)$ définie pour tout entier naturel $n \geqslant 1$ par:
$v_n = 9 - 6 \times \e^{-0,19\times(n - 1)}$.

Le terme $v_n$ est une estimation du nombre de questions, en centaines, présentes le $n$-ième mois sur la FAQ.

\medskip

\begin{enumerate}
\item 
$v_1=9-6\e^{0}=3$ et $v_2=9-6\e^{-0,19}\approx 4,04$

\item On résout l'inéquation $v_n > 8,5$.

$v_n > 8,5
\iff
9 - 6 \times \e^{-0,19\times(n - 1)} > 8,5
\iff 
0,5 > 6 \times \e^{-0,19\times(n - 1)}
$\\

$\phantom{v_n > 8,5}\iff\dfrac{0,5}{6} > \e^{-0,19\times(n - 1)}$ et par croissance de la fonction logarithme népérien\\
$\phantom{v_n > 8,5}
\iff
\ln \left (\dfrac{0,5}{6}\right ) > -0,19\times (n-1)
\iff
-\dfrac{\ln \left (\frac{0,5}{6}\right )}{0,19} < n - 1 \iff n > 1 -\dfrac{\ln \left (\frac{0,5}{6}\right )}{0,19}$

Or $-\dfrac{\ln \left (\frac{0,5}{6}\right )}{0,19}\approx -13,08$ et 1 + 13,08 = 14,08 donc la plus petite valeur telle que $v_n > 8,5$ est $n = 15$.

\end{enumerate}

\bigskip

\textbf{Partie C : Comparaison des deux modèles}

\medskip

\begin{enumerate}
\item L'entreprise considère qu'elle doit modifier la présentation de son site lorsque plus de 850 questions sont présentes sur la FAQ. 

\begin{list}{\textbullet}{}
\item Selon le 1\ier{} modèle, il y a plus de 850 questions sur la FAQ quand $u_n>8,5$, c'est-à-dire le 9\up{e} mois.
\item Selon le 2\up{e} modèle, il y a plus de 850 questions sur la FAQ quand $v_n>8,5$, c'est-à-dire le 15\up{e} mois.
\end{list}

C'est la première modélisation qui conduit à procéder le plus tôt à cette modification.


\item Pour savoir pour quelle modélisation il y a le plus grand nombre de questions sur la FAQ à long terme, on cherche les limites des suites $(u_n)$ et $(v_n)$.

\begin{list}{\textbullet}{}
\item $u_n=13 - \dfrac{100}{9}\times 0,9^n$

$-1<0,9 <1$ donc $\ds\lim_{n\to +\infty} 0,9^n=0$ donc $\ds\lim_{n\to +\infty} u_n=13$.

À long terme il y aura \np{1300} questions pour la 1\up{e} modélisation..
\item $v_n = 9 - 6 \times \e^{-0,19\times(n - 1)}$

$\ds\lim_{n\to +\infty} -0,19\times (n-1) = -\infty$ et $\ds\lim_{X\to -\infty} \e^{X}=0$ donc $\ds\lim_{n\to +\infty} \e^{-0,19\times(n - 1)}=0$; on en déduit que $\ds\lim_{n\to +\infty} v_n=9$.

À long terme il y aura \np{900} questions pour la 2\up{e} modélisation..
\end{list}

C'est donc pour la 1\up{re} modélisation qu'il y aura le plus de questions à long terme.

\end{enumerate}

\bigskip
